\documentclass[11pt]{article}
% Shared preamble for the rigor documents (Lamport-style structured proofs).  \input this file after \documentclass.
\usepackage[T1]{fontenc}\usepackage{lmodern}\usepackage{microtype}
\usepackage[margin=1in]{geometry}
\usepackage{amsmath,amssymb,amsthm,mathtools}
\usepackage{xcolor}\usepackage{enumitem}\usepackage{booktabs}\usepackage{graphicx}\usepackage{hyperref}
\usepackage[most]{tcolorbox}
\definecolor{navy}{HTML}{17365D}\definecolor{teal}{HTML}{117A8B}\definecolor{warn}{HTML}{A33A2B}\definecolor{okgreen}{HTML}{2E7D4F}
\hypersetup{colorlinks=true,linkcolor=navy,citecolor=teal,urlcolor=teal}
\theoremstyle{plain}
\newtheorem{theorem}{Theorem}[section]\newtheorem{lemma}[theorem]{Lemma}\newtheorem{proposition}[theorem]{Proposition}\newtheorem{corollary}[theorem]{Corollary}
\theoremstyle{definition}\newtheorem{definition}[theorem]{Definition}\newtheorem{remark}[theorem]{Remark}\newtheorem{claim}[theorem]{Claim}\newtheorem{issue}[theorem]{Issue}
% ---------- Lamport structured proofs ----------
% \lstep{level}{number}{assertion}   -> indented "<level>number. assertion"
% \lproof                             -> "PROOF:" line (start of the sub-proof of the preceding step)
% \lby{justification}                 -> "BY ..." leaf justification for the preceding step
% \lqed{level}{number}                -> QED step of a sub-proof
% keywords: \lassume \lprove \lsuffices \lcase \ldefine \lpick \lnew
\newlength{\lindent}\setlength{\lindent}{1.7em}
\newcommand{\lstep}[3]{\par\smallskip\noindent\hspace*{\dimexpr\numexpr#1-1\relax\lindent\relax}{\color{navy}$\langle#1\rangle#2$.}\ #3}
\newcommand{\lproof}{\par\noindent\hspace*{\lindent}{\scshape Proof:}}
\newcommand{\lby}[1]{\par\noindent\hspace*{\lindent}{\scshape By}\ #1.}
\newcommand{\lqed}[2]{\par\smallskip\noindent\hspace*{\dimexpr\numexpr#1-1\relax\lindent\relax}{\color{navy}$\langle#1\rangle#2$.}\ {\scshape Q.E.D.}}
\newcommand{\lassume}{{\scshape Assume}\ }\newcommand{\lprove}{{\scshape Prove}\ }\newcommand{\lsuffices}{{\scshape Suffices}\ }
\newcommand{\lcase}{{\scshape Case}\ }\newcommand{\ldefine}{{\scshape Define}\ }\newcommand{\lpick}{{\scshape Pick}\ }\newcommand{\lnew}{{\scshape New}\ }
% ---------- verdict boxes ----------
\newtcolorbox{verdictok}{colback=okgreen!8,colframe=okgreen,title={\bfseries Verdict: verified},fonttitle=\small}
\newtcolorbox{verdictgap}{colback=orange!10,colframe=orange!80!black,title={\bfseries Verdict: gap / caveat},fonttitle=\small}
\newtcolorbox{verdictbreak}{colback=warn!10,colframe=warn,title={\bfseries Verdict: proof-breaking issue},fonttitle=\small}
\newtcolorbox{citedbox}[1]{colback=navy!5,colframe=navy!60,title={\bfseries Cited result: #1},fonttitle=\small,breakable}
\setlength{\parindent}{0pt}\setlength{\parskip}{4pt}\allowdisplaybreaks
\newcommand{\cA}{\mathcal A}\newcommand{\cF}{\mathcal F}\newcommand{\cM}{\mathfrak M}\newcommand{\cN}{\mathfrak N}

% extra calligraphic shortcuts (\providecommand: no clash if a document defines its own)
\providecommand{\cB}{\mathcal B}\providecommand{\cC}{\mathcal C}\providecommand{\cD}{\mathcal D}\providecommand{\cE}{\mathcal E}\providecommand{\cG}{\mathcal G}\providecommand{\cH}{\mathcal H}\providecommand{\cI}{\mathcal I}\providecommand{\cJ}{\mathcal J}\providecommand{\cK}{\mathcal K}\providecommand{\cL}{\mathcal L}\providecommand{\cO}{\mathcal O}\providecommand{\cP}{\mathcal P}\providecommand{\cQ}{\mathcal Q}\providecommand{\cR}{\mathcal R}\providecommand{\cS}{\mathcal S}\providecommand{\cT}{\mathcal T}\providecommand{\cU}{\mathcal U}\providecommand{\cV}{\mathcal V}\providecommand{\cW}{\mathcal W}\providecommand{\cX}{\mathcal X}\providecommand{\cY}{\mathcal Y}\providecommand{\cZ}{\mathcal Z}
\newcommand{\Fid}{\operatorname{F}}\newcommand{\Tr}{\operatorname{Tr}}\newcommand{\eps}{\varepsilon}\newcommand{\dd}{\,\mathrm d}

\newcommand{\hh}{\mathfrak h}
\newcommand{\Hrm}{\mathsf H}
\newcommand{\Erec}{E_{\mathrm{opt}}}
\newcommand{\Krec}{\mathcal K_{\mathrm{rec}}}
\title{Second-order optimality of the Petz recovery map\\
\large Phase~3, track S8 (problem O8): formalisation, the convex quasi-free
programme, numerics, and a lower bound}
\author{Agent S8}
\date{2026-09-09}
\begin{document}
\maketitle

\begin{tcolorbox}[colback=navy!5,colframe=navy!60,title={\bfseries Summary},fonttitle=\small]
The optimal recovery error \(\Erec\) is formalised as an infimum of the convex
function \(-\log\Fid(\omega,\cdot)\) over the convex set \(\Krec\) of recovered
states (Prop.~\ref{prop:convex}).  For the free fermion and gauge-invariant
quasi-free channels the problem is shown to be an \emph{exactly convex}
programme in the pair \((X,Z)=(\text{one-particle map},\text{output symbol})\),
with feasible set \(XQ_BX^*\le Z\le 1-X(1-Q_B)X^*\) (Thm.~\ref{thm:convex-set});
the complete-positivity condition is verified against an explicit Stinespring
dilation on three modes.  Numerically, on the half-filled hopping chain, the
optimum beats the Petz map by a factor \(\approx4\) and its coefficient agrees
with \emph{twice} the single-chirality geometric-compression value
\(f_2z^2\), i.e.\ with the two-chirality spatial compression: to the accuracy
reached, \textbf{the Petz map is not second-order optimal, and the geometric
compression is} (or is within a few percent of the optimum).  A lower bound
valid for \emph{all} channels is obtained by restricting to \(\cA(A)\vee\cA(C)\)
and to a single-mode measurement; it is nonzero at order \(z^2\) but is two
orders of magnitude below the optimum.  Gauge twirling reduces the problem to
gauge-covariant channels but \emph{not} to quasi-free ones; this is stated
honestly as the remaining gap.
\end{tcolorbox}

\tableofcontents
\clearpage

\section{Setting, conventions, and the optimal recovery error}\label{sec:setting}

We keep the conventions of Notes~4 and~5 (\texttt{continuum\_petz\_free\_fermion.tex},
\texttt{fidelity\_recovered\_quasifree.tex}).  Thus \(A=(-a,0)\), \(B=(0,b)\),
\(C=(b,L)\), \(L=b+c\), \(D:=BC=(0,L)\), \(I:=ABC=(-a,L)\); \(\omega\) is the
vacuum state of the net \(X\mapsto\cA(X)\) (for the free chiral fermion,
\(\cA(X)=\cF_r(X)\)); and
\begin{equation}\label{eq:zdef}
 z:=\frac{ac}{b(a+b+c)}=\frac1\eta-1,\qquad
 \eta_{\mathrm V}=\frac{ac}{(a+b)(b+c)}=\frac{z}{1+z},
\end{equation}
is the cross-ratio variable of Note~5, \(\zeta=as/(a+L)\) the parameter of the
geometric family, \(\zeta=2z\) for the ordinary Petz map and \(\zeta=z\) for the
zero-collar geometric compression.

\begin{definition}[Admissible recovery channel]\label{def:admissible}
A \emph{recovery channel} is a map \(\beta:\cA(BC)\to\cA(B)\) that is
(i) linear, (ii) completely positive, (iii) unital, \(\beta(1)=1\),
(iv) normal (\(\sigma\)-weakly continuous), and, for fermionic nets,
(v) even, i.e.\ \(\beta\circ\alpha^{\mathrm{gr}}=\alpha^{\mathrm{gr}}\circ\beta\)
for the grading automorphism \(\alpha^{\mathrm{gr}}\).  We write \(\mathcal C\)
for the set of recovery channels.  Condition (v) guarantees that
\(\alpha_\beta:=\mathrm{id}_{\cA(A)}\otimes\beta\) is well defined and
completely positive on the graded tensor product
\(\cA(A)\,\widehat\otimes\,\cA(BC)\supseteq\cA(ABC)\); for the split inclusions
used in Notes~4--5 this is the tensor product of the two type-I factors
interpolating the collar (Note~4, Sec.~2).
\end{definition}

\begin{definition}[Recovered state, optimal recovery error]\label{def:Eopt}
For \(\beta\in\mathcal C\) the \emph{recovered state} is the normal state on
\(\cA(ABC)\)
\begin{equation}\label{eq:rec-state}
 \omega^\beta:=\bigl(\omega|_{\cA(AB)}\bigr)\circ\alpha_\beta,\qquad
 \omega^\beta(x\,y)=\omega\bigl(x\,\beta(y)\bigr)\quad(x\in\cA(A),\,y\in\cA(BC)),
\end{equation}
and the \emph{optimal recovery error} is
\begin{equation}\label{eq:Eopt}
 \boxed{\;\Erec:=\inf_{\beta\in\mathcal C}\bigl[-\log\Fid(\omega,\omega^\beta)\bigr]\;}
\end{equation}
with \(\Fid\) the Uhlmann root fidelity (Note~5, Sec.~1).  We also write
\(\Krec:=\{\omega^\beta:\beta\in\mathcal C\}\) for the set of recovered states.
\end{definition}

\begin{proposition}[Elementary properties]\label{prop:convex}
\begin{enumerate}[leftmargin=1.6em,itemsep=1pt]
\item \(\mathcal C\) is convex and closed under composition with unital channels
 on either side; \(\beta\mapsto\omega^\beta\) is affine, hence \(\Krec\) is convex.
\item \(-\log\Fid(\omega,\cdot)\) is convex on the state space; hence
 \eqref{eq:Eopt} is the infimum of a convex function on a convex set, and every
 local minimum is global.
\item \(\Erec\ge0\), and \(\Erec=0\) iff \(\omega\) is a quantum Markov state for
 \(A\!-\!B\!-\!C\).
\item Every rotated Petz map and the zero-collar geometric compression belong to
 \(\mathcal C\); consequently, per chirality,
 \(\Erec\le\Phi(z)=f_2z^2+O(z^3)\) with \(f_2=c/(12\pi^2)\), whereas the ordinary
 Petz map gives \(\Phi(2z)=4f_2z^2+O(z^3)\).  For a two-dimensional CFT the
 spatial compression acts identically on both chiralities and gives
 \(\Erec\le2f_2z^2+O(z^3)\) against the Petz value \(8f_2z^2+O(z^3)\).
\end{enumerate}
\end{proposition}

\begin{proof}
\lstep{1}{1}{\(\mathcal C\) is convex and \(\beta\mapsto\omega^\beta\) is affine.}
\lby{complete positivity, unitality, normality and evenness are each preserved
by convex combinations (for complete positivity: \(\bigl[(\lambda\beta_1+(1-\lambda)\beta_2)(x_i^*x_j)\bigr]_{ij}
=\lambda[\beta_1(x_i^*x_j)]_{ij}+(1-\lambda)[\beta_2(x_i^*x_j)]_{ij}\ge0\) as a sum
of positive matrices); and \eqref{eq:rec-state} is linear in \(\beta\)}
\lstep{1}{2}{\(\Krec\) is convex.}
\lby{$\langle1\rangle1$ and the fact that the image of a convex set under an
affine map is convex}
\lstep{1}{3}{\(\varphi\mapsto-\log\Fid(\omega,\varphi)\) is convex.}
\lproof
\lstep{2}{1}{\(\Fid\) is jointly concave in its two arguments.}
\lby{Uhlmann 1976 (\emph{Rep.\ Math.\ Phys.}\ \textbf{9}, 273), joint concavity of
the transition probability; in finite dimensions Nielsen--Chuang Thm.~9.7}
\lstep{2}{2}{\(t\mapsto-\log t\) is convex and nonincreasing on \((0,1]\).}
\lby{elementary}
\lqed{2}{3}
\lby{the composition of a nonincreasing convex function with a concave function
is convex, applied to $\langle2\rangle1$--$\langle2\rangle2$; if
\(\Fid(\omega,\varphi)=0\) the value is \(+\infty\) and convexity persists}
\lstep{1}{4}{\(\Erec\ge0\), with equality iff \(\omega\) is Markov.}
\lproof
\lstep{2}{1}{\(\Fid\le1\) always, so \(-\log\Fid\ge0\).}
\lby{Cauchy--Schwarz for vector representatives (Note~5, Sec.~1)}
\lstep{2}{2}{If \(\Erec=0\) then there are \(\beta_n\) with
\(\Fid(\omega,\omega^{\beta_n})\to1\), hence \(\omega^{\beta_n}\to\omega\) in norm;
conversely a Markov state admits an exact recovery.}
\lby{the Fuchs--van~de~Graaf inequality
\(\tfrac12\|\varphi-\psi\|_1\le\sqrt{1-\Fid^2}\) and, for the converse, the Petz
map of a Markov state is an exact recovery (Petz 1986, 1988; Note~3)}
\lqed{2}{3}
\lstep{1}{5}{The geometric maps are admissible and give the stated values.}
\lproof
\lstep{2}{1}{For \(s\ge\lambda:=c/b\) the map \(\gamma_s\) of Note~4 is a normal
unital \(*\)-endomorphism \(\cA(D)\to\cA(J_s\cap(0,\infty))\subseteq\cA(B)\), hence
lies in \(\mathcal C\).}
\lby{Note~4, Sec.~2: \(\gamma_s=\operatorname{Ad}U(h_s)\) restricted to
\(\cA(D)\), \(h_s(0,L)=(0,L/(1+s))\) and \(L/(1+s)\le b\iff s\ge\lambda\);
a unital \(*\)-homomorphism is completely positive and even}
\lstep{2}{2}{\(-\log\Fid(\omega,\omega\circ\overline\gamma_s)=\Phi(\zeta)\),
\(\zeta=as/(a+L)\), and \(\Phi\) is nondecreasing.}
\lby{Note~5, Lemma~1.1 and Cor.~5.1, monotonicity proved in
\texttt{rigor/check\_note5\_theory.tex}}
\lstep{2}{3}{The constraint \(s\ge\lambda\) is exactly \(\zeta\ge z\), and the
infimum over the geometric family is attained at \(s=\lambda\), value
\(\Phi(z)=f_2z^2+O(z^3)\).}
\lby{$\langle2\rangle1$--$\langle2\rangle2$, \(a\lambda/(a+L)=ac/(b(a+b+c))=z\)
by \eqref{eq:zdef}, and Note~7 for \(f_2=c/(12\pi^2)\)}
\lqed{2}{4}
\lby{the ordinary Petz map is \(s=2\lambda\), \(\zeta=2z\), and \(\Phi\) is
quadratic to leading order}
\lqed{1}{6}
\end{proof}

\begin{remark}[The constraint is active]\label{rem:active}
Step \(\langle2\rangle3\) is the structural reason why the geometric compression
is the best member of the rotated family: \emph{the range constraint
\(\beta(\cA(BC))\subseteq\cA(B)\) is active at the optimum}.  Any second-order
optimality statement must therefore be a constrained one; an unconstrained
stationarity argument cannot prove it.
\end{remark}



\section{Gauge-invariant quasi-free channels: parametrisation and convexity}
\label{sec:qf}

Let \(\hh_A,\hh_B,\hh_C\) be the one-particle spaces of \(A,B,C\),
\(\hh_D=\hh_B\oplus\hh_C\), \(\hh_I=\hh_A\oplus\hh_D\).  A gauge-invariant
quasi-free state \(\omega_Q\) is fixed by its symbol,
\(\omega_Q(a^*(f)a(g))=\langle g,Qf\rangle\); we use throughout the matrix
convention \(Q_{ij}=\omega(c_j^*c_i)\), for which a number-conserving Bogoliubov
map \(c\mapsto Wc\) acts as \(Q\mapsto WQW^*\).  The vacuum symbol on \(I\) is
\(Q=E_IP_+E_I\) (continuum) resp.\ the sine kernel
\(Q_{ij}=\frac{\sin(\pi(i-j)/2)}{\pi(i-j)}\) (half-filled hopping chain), with
blocks \(Q_{AA},Q_{AB},Q_{AC},Q_{BB},Q_{BC},Q_{CC}\).

\begin{definition}[Quasi-free recovery channel]\label{def:qf-channel}
A recovery channel \(\beta\) is \emph{gauge-invariant quasi-free} if its
Schr\"odinger picture \(\cN=\beta_*:\mathcal S(\cA(B))\to\mathcal S(\cA(BC))\)
maps quasi-free states to quasi-free states by an affine rule on symbols,
\begin{equation}\label{eq:channel-rule}
 Q^{\mathrm{out}}=X\,Q^{\mathrm{in}}X^*+Y,\qquad X:\hh_B\to\hh_D,\quad Y=Y^*\ \text{on}\ \hh_D .
\end{equation}
\end{definition}

\begin{proposition}[Complete positivity]\label{prop:cp}
The rule \eqref{eq:channel-rule} is the symbol rule of a completely positive,
trace-preserving, even map \(\cN\) if and only if
\begin{equation}\label{eq:cp}
 \boxed{\;Y\ge0\quad\text{and}\quad XX^*+Y\le1 .\;}
\end{equation}
Under \eqref{eq:cp} a Stinespring dilation with a \emph{Gaussian} environment
exists: with \(P:=1-XX^*\ge Y\), \(V_E:=P^{1/2}\), \(Q_E:=P^{-1/2}YP^{-1/2}\)
(on \(\operatorname{ran}P\), \(0\le Q_E\le1\)) and any unitary \(W\) on
\(\hh_B\oplus\hh_E\) whose first \(\dim\hh_D\) rows are \((X\ V_E)\), one has
\(\cN(\rho)=\Tr_{\hh_{\mathrm{disc}}}\bigl[\Gamma(W)(\rho\otimes\rho_{Q_E})\Gamma(W)^*\bigr]\).
\end{proposition}

\begin{proof}
\lstep{1}{1}{\lassume{} \eqref{eq:channel-rule} comes from a positive
trace-preserving map. \lprove{} \eqref{eq:cp}.}
\lby{the quasi-free states with \(Q^{\mathrm{in}}=0\) and \(Q^{\mathrm{in}}=1\)
are states (the Fock vacuum and the filled state); their images must have
symbols in \([0,1]\), giving \(Y\ge0\) and \(XX^*+Y\le1\)}
\lstep{1}{2}{\lassume{} \eqref{eq:cp}. \lprove{} the dilation exists.}
\lproof
\lstep{2}{1}{\(0\le Y\le P:=1-XX^*\), and \(Q_E:=P^{-1/2}YP^{-1/2}\) satisfies
\(0\le Q_E\le1\) on \(\operatorname{ran}P\) (set \(Q_E=0\) on \(\ker P\), where
\(Y=0\) because \(0\le Y\le P\)).}
\lby{\eqref{eq:cp} and operator monotonicity of \(t\mapsto t\) under congruence:
\(Y\le P\Rightarrow P^{-1/2}YP^{-1/2}\le1\)}
\lstep{2}{2}{\(V:=(X\ \ P^{1/2}):\hh_B\oplus\hh_E\to\hh_D\) satisfies
\(VV^*=XX^*+P=1_{\hh_D}\), so \(V\) is a coisometry and extends to a unitary
\(W\) on \(\hh_B\oplus\hh_E\) with \(\dim\hh_E=\dim\hh_D\) (complete the rows of
\(V\) to an orthonormal basis).}
\lby{linear algebra; \(\dim(\hh_B\oplus\hh_E)\ge\dim\hh_D\)}
\lstep{2}{3}{The map \(\rho\mapsto\Tr_{\mathrm{disc}}[\Gamma(W)(\rho\otimes\rho_{Q_E})\Gamma(W)^*]\)
is completely positive, trace preserving and even, and its symbol rule is
\(Q^{\mathrm{out}}=V(Q^{\mathrm{in}}\oplus Q_E)V^*=XQ^{\mathrm{in}}X^*+P^{1/2}Q_EP^{1/2}
=XQ^{\mathrm{in}}X^*+Y\).}
\lby{\(\Gamma(W)\) is a unitary implementer of an even Bogoliubov
transformation, the partial trace over a set of modes is CP, unital in the
Heisenberg picture and even, and both operations act on two-point functions by
compression with \(V\)}
\lqed{2}{4}
\lqed{1}{3}
\end{proof}

\begin{verdictok}
Both directions of Prop.~\ref{prop:cp} were checked numerically
(\texttt{numerics/optimality/cp\_check.py}): random \((X,Y)\) obeying
\eqref{eq:cp} on \(\dim\hh_B=1,\dim\hh_D=2\) were dilated as above, the full
\(2^4\)-dimensional Fock-space channel built explicitly, and the resulting
output two-point function compared with \eqref{eq:channel-rule}; agreement to
\(10^{-14}\).  Violating either inequality produced an output symbol outside
\([0,1]\) for an extreme input, hence a non-positive map.
\end{verdictok}

\begin{lemma}[Recovered symbol]\label{lem:rec-symbol}
Let \(\beta\) be gauge-invariant quasi-free with data \((X,Y)\) and put
\(Z:=XQ_{BB}X^*+Y\).  Then \(\omega^\beta\) is the gauge-invariant quasi-free
state on \(\cA(I)\) with symbol
\begin{equation}\label{eq:Qrec}
 Q^{\mathrm{rec}}=\begin{pmatrix}Q_{AA}&Q_{AB}X^*\\ XQ_{BA}&Z\end{pmatrix}
 \quad\text{on }\hh_A\oplus\hh_D .
\end{equation}
Writing \(X^*=(S\ \ T)\) with \(S:\hh_B\to\hh_B\), \(T:\hh_C\to\hh_B\), the
recovered \(A\!-\!C\) block is \(Q^{\mathrm{rec}}_{AC}=Q_{AB}T\): \emph{the
recovered \(A\!-\!C\) correlations are vacuum \(A\!-\!B\) correlations evaluated
on the image of \(T\).}
\end{lemma}

\begin{proof}
\lstep{1}{1}{\(\alpha_\beta=\mathrm{id}_A\otimes\beta\) acts on the vacuum
\(AB\) state by applying \(\cN\) to the \(B\) tensor factor.}
\lby{Definition~\ref{def:Eopt} and evenness (Def.~\ref{def:admissible}(v)),
which makes \(\mathrm{id}\otimes\cN\) well defined on the graded tensor product}
\lstep{1}{2}{In the dilation of Prop.~\ref{prop:cp} the global state is
\(\rho_{AB}\otimes\rho_{Q_E}\), a quasi-free state with symbol
\(Q_{AB\text{-block}}\oplus Q_E\); \(\Gamma(1_A\oplus W)\) is Bogoliubov and the
partial trace is a compression, so the image is quasi-free with symbol
\((1_A\oplus V)(Q^{(AB)}\oplus Q_E)(1_A\oplus V)^*\).}
\lby{$\langle1\rangle1$, Prop.~\ref{prop:cp}, and the fact that the vacuum
\(AB\) state is quasi-free with symbol \(\bigl(\begin{smallmatrix}Q_{AA}&Q_{AB}\\Q_{BA}&Q_{BB}\end{smallmatrix}\bigr)\)}
\lstep{1}{3}{That symbol equals \eqref{eq:Qrec}.}
\lby{block multiplication: the \(AA\) block is \(Q_{AA}\), the \(AD\) block is
\(Q_{AB}V^*|_{\hh_B\text{-part}}=Q_{AB}X^*\) (the environment is uncorrelated
with \(A\)), and the \(DD\) block is \(XQ_{BB}X^*+P^{1/2}Q_EP^{1/2}=Z\)}
\lqed{1}{4}
\end{proof}



\section{The second-order problem is an exactly convex programme}\label{sec:convex}

\begin{theorem}[Convex reparametrisation]\label{thm:convex-set}
Put \(R_+:=Q_{BB}\), \(R_-:=1-Q_{BB}\) (both \(\ge0\)) and
\begin{equation}\label{eq:feasible}
 \mathcal F:=\bigl\{(X,Z):\ X:\hh_B\to\hh_D,\ Z=Z^*\ \text{on}\ \hh_D,\
  XR_+X^*\ \le\ Z\ \le\ 1-XR_-X^*\bigr\}.
\end{equation}
Then
\begin{enumerate}[leftmargin=1.6em,itemsep=1pt]
\item \(\mathcal F\) is exactly the set of pairs \((X,Z)\) with
 \(Z=XQ_{BB}X^*+Y\) for some \((X,Y)\) satisfying the complete-positivity
 condition \eqref{eq:cp};
\item \(\mathcal F\) is convex;
\item \(Q^{\mathrm{rec}}\) of \eqref{eq:Qrec} is an affine function of \((X,Z)\);
\item consequently, for any convex functional \(\mathcal G\) of the symbol, the
 problem \(\min_{(X,Z)\in\mathcal F}\mathcal G(Q^{\mathrm{rec}}(X,Z))\) is a
 convex programme whose constraints are two linear matrix inequalities in the
 Schur-complement variables; every local minimum is global.
\end{enumerate}
\end{theorem}

\begin{proof}
\lstep{1}{1}{(1) holds.}
\lby{substituting \(Y=Z-XR_+X^*\): \(Y\ge0\iff Z\ge XR_+X^*\), and
\(XX^*+Y\le1\iff Z\le1-XX^*+XR_+X^*=1-XR_-X^*\)}
\lstep{1}{2}{For \(R\ge0\) the map \(X\mapsto XRX^*\) is operator convex.}
\lby{for \(X_t=(1-t)X_0+tX_1\), \(t\in[0,1]\), the identity
\((1-t)X_0RX_0^*+tX_1RX_1^*-X_tRX_t^*=t(1-t)(X_0-X_1)R(X_0-X_1)^*\ge0\)}
\lstep{1}{3}{\(\mathcal F\) is convex.}
\lby{$\langle1\rangle2$: \(\{(X,Z):Z-XR_+X^*\ge0\}\) and
\(\{(X,Z):1-XR_-X^*-Z\ge0\}\) are the sublevel sets of two operator-convex
functions of \((X,Z)\) (each affine in \(Z\)), hence convex; the intersection of
convex sets is convex}
\lstep{1}{4}{(3) and (4).}
\lby{\eqref{eq:Qrec} is affine in \((X,Z)\) by inspection; the composition of a
convex functional with an affine map is convex, and a convex function on a
convex set has no non-global local minima}
\lqed{1}{5}
\end{proof}

\begin{remark}
The reparametrisation is essential.  In the natural variables \((X,Y)\) the
objective is a \emph{quartic} function of \(X\) (because \(Q^{\mathrm{rec}}_{DD}\)
contains \(XQ_{BB}X^*\)) and the problem is not convex; in \((X,Z)\) the
nonconvexity is absorbed into the constraints, which become operator convex.
The set \(\mathcal F\) is semidefinite representable: writing the two constraints
as \(\bigl(\begin{smallmatrix}Z&XR_+^{1/2}\\ R_+^{1/2}X^*&1\end{smallmatrix}\bigr)\ge0\)
and \(\bigl(\begin{smallmatrix}1-Z&XR_-^{1/2}\\ R_-^{1/2}X^*&1\end{smallmatrix}\bigr)\ge0\)
(Schur complements) shows that the second-order optimality problem is a
semidefinite programme.
\end{remark}

\subsection{The second-order objective and the order of limits}\label{sec:order}

By Note~5, Prop.~2.3, for a fixed reference symbol \(Q\) with \(0<Q<1\),
\begin{equation}\label{eq:gQ}
 -\log\Fid(\omega_Q,\omega_{Q-\delta})=g_Q(\delta)+O(\|\delta\|^3),\qquad
 g_Q(\delta):=\frac14\sum_{i,k}w_{ik}\,|\widetilde\delta_{ik}|^2,\quad
 w_{ik}=\frac{1}{q_i(1-q_k)+q_k(1-q_i)},
\end{equation}
\(\widetilde\delta=U^*\delta U\) in an eigenbasis of \(Q\), and
\(w_{ik}=1+\cosh\frac{\kappa_i+\kappa_k}2/\cosh\frac{\kappa_i-\kappa_k}2\) in the
modular variables \(q=(1+e^\kappa)^{-1}\).  \(g_Q\) is a positive quadratic form,
so Theorem~\ref{thm:convex-set}(4) applies:

\begin{corollary}[Second-order optimality is a convex QP]\label{cor:convexQP}
The second-order optimal recovery error over gauge-invariant quasi-free channels,
\begin{equation}\label{eq:secondorderproblem}
 \Erec^{(2)}:=\min_{(X,Z)\in\mathcal F}\ g_Q\bigl(Q-Q^{\mathrm{rec}}(X,Z)\bigr),
\end{equation}
is the minimum of a positive quadratic form over a convex, semidefinite
representable set.  The minimum is attained (the form is continuous and, after
restricting to a bounded slice \(\|X\|\le1\), \(\mathcal F\) is compact).
\end{corollary}

\begin{verdictgap}[Order of limits]
\(g_Q\) is the \emph{second-order} term of \(-\log\Fid\) only when the defect is
small \emph{relative to the extreme eigenvalues of \(Q\)}: the weight
\(w_{ik}\approx e^{\min(|\kappa_i|,|\kappa_k|)}\) amplifies defect matrix elements
between two nearly empty (or two nearly filled) modular modes, so
\eqref{eq:gQ} requires \(\|\delta\|\ll e^{-\kappa_{\max}}\).  On a chain of
\(n=13\) sites the Petz defect at \(z=0.15\) has
\(-\log\Fid=1.303\times10^{-3}\) but \(g_Q(\delta)=3.61\times10^{-2}\), a factor
\(28\): at accessible \(z\) the \emph{unrestricted} form \eqref{eq:gQ} is not yet
the second-order approximation.  Two consequences, both handled below.
(i) Minimising \(g_Q\) restricted to a modular window \(|\kappa|\le\kappa_c\)
(as in Note~5) is a valid lower bound for \eqref{eq:secondorderproblem} but
produces a \emph{useless minimiser}: the optimiser dumps the defect into the
discarded modes; on \((L_A,L_B,L_C)=(4,6,3)\), \(\kappa_c=10\), the windowed
minimiser has an exact \(-\log\Fid=1.145\), i.e.\ \(880\) times \emph{worse} than
Petz.  (ii) The cure used here is to \emph{cap} rather than discard:
\(w\mapsto\min(w,1+\cosh\kappa_c)\).  Capping preserves convexity and still lower
bounds \(g_Q\), but no mode is free, so the minimiser stays sensible; the capped
programme is then used only as a \emph{generator of candidate channels}, and the
arbiter is the exact high-precision \(-\log\Fid\) of the recovered state.  Since
a candidate channel is a certificate (it is admissible by construction, and its
error is computed exactly), all \emph{upper} bounds reported below are rigorous
within the lattice model regardless of how the channel was found.
\end{verdictgap}



\section{Lower bounds: what survives and what does not}\label{sec:lower}

\subsection{The restriction bound and the \(A\!-\!C\) route}

\begin{lemma}[Restriction]\label{lem:restriction}
For any von Neumann subalgebra \(\cN\subseteq\cA(ABC)\),
\(\Fid(\omega,\omega^\beta)\le\Fid(\omega|_\cN,\omega^\beta|_\cN)\); in particular
\begin{equation}\label{eq:AC-bound}
 \Erec\ \ge\ E_{AC}:=-\log\ \sup_{\cN_{B\to C}}\
 \Fid\Bigl(\omega|_{\cA(A)\vee\cA(C)},\ (\mathrm{id}_A\otimes\cN_{B\to C})\bigl[\omega|_{\cA(A)\vee\cA(B)}\bigr]\Bigr),
\end{equation}
the supremum running over all channels from \(B\) to \(C\).
\end{lemma}

\begin{proof}
\lstep{1}{1}{The inclusion \(\iota:\cN\hookrightarrow\cA(ABC)\) is unital and
completely positive, and \(\Fid\) is monotone nondecreasing under the pull-back
of states along a unital CP map.}
\lby{Alberti--Uhlmann 2002, Thm.~2 (the variational formula
\(\Fid(\varphi,\psi)=\inf_{x>0}\frac12(\varphi(x)+\psi(x^{-1}))\) is manifestly
monotone); Note~5, Sec.~1}
\lstep{1}{2}{For \(x\in\cA(A)\), \(y\in\cA(C)\subseteq\cA(BC)\),
\(\omega^\beta(xy)=\omega(x\,\beta(y))\) with \(\beta(y)\in\cA(B)\); hence
\(\omega^\beta|_{\cA(A)\vee\cA(C)}\) is the image of \(\omega|_{\cA(A)\vee\cA(B)}\)
under \(\mathrm{id}_A\otimes\cN\), \(\cN=(\beta|_{\cA(C)})_*\).}
\lby{Definition~\ref{def:Eopt}; \(\beta|_{\cA(C)}:\cA(C)\to\cA(B)\) is unital CP}
\lqed{1}{3}
\lby{$\langle1\rangle1$ with \(\cN=\cA(A)\vee\cA(C)\) and $\langle1\rangle2$}
\end{proof}

\begin{proposition}[The \(A\!-\!C\) bound vanishes]\label{prop:AC-zero}
Suppose there is \(T:\hh_C\to\hh_B\) with
\begin{equation}\label{eq:Tcond}
 Q_{AB}T=Q_{AC},\qquad T^*T\le1,\qquad
 T^*Q_{BB}T\ \le\ Q_{CC}\ \le\ 1-T^*(1-Q_{BB})T .
\end{equation}
Then there is a gauge-invariant quasi-free \(\beta\in\mathcal C\) with
\(\omega^\beta|_{\cA(A)\vee\cA(C)}=\omega|_{\cA(A)\vee\cA(C)}\) exactly, so
\(E_{AC}=0\) and \eqref{eq:AC-bound} is empty.
\end{proposition}

\begin{proof}
\lstep{1}{1}{\ldefine{} \(X^*:=(0\ \ T):\hh_B\oplus\hh_C\to\hh_B\) and
\(Z:=\tfrac12 1_{\hh_B}\oplus Q_{CC}\).}
\lstep{1}{2}{\((X,Z)\in\mathcal F\).}
\lby{\(XR_\pm X^*=0\oplus T^*R_\pm T\) is block diagonal, so
\(Z-XR_+X^*=\tfrac12 1\oplus(Q_{CC}-T^*Q_{BB}T)\ge0\) and
\(1-XR_-X^*-Z=\tfrac12 1\oplus(1-T^*(1-Q_{BB})T-Q_{CC})\ge0\) by \eqref{eq:Tcond}}
\lstep{1}{3}{The recovered symbol restricted to \(\hh_A\oplus\hh_C\) is
\(\bigl(\begin{smallmatrix}Q_{AA}&Q_{AB}T\\ T^*Q_{BA}&Q_{CC}\end{smallmatrix}\bigr)
=\bigl(\begin{smallmatrix}Q_{AA}&Q_{AC}\\ Q_{CA}&Q_{CC}\end{smallmatrix}\bigr)\).}
\lby{Lemma~\ref{lem:rec-symbol} and \eqref{eq:Tcond}}
\lqed{1}{4}
\lby{two gauge-invariant quasi-free states with the same symbol on
\(\hh_A\oplus\hh_C\) coincide on \(\cA(A)\vee\cA(C)\), and by
Prop.~\ref{prop:cp} \((X,Z)\) comes from an admissible channel}
\end{proof}

\begin{verdictbreak}[\;The \(A\!-\!C\) route of the O8 plan does not work]
Hypothesis \eqref{eq:Tcond} was checked on the hopping chain with
\(T:=Q_{AB}^{+}Q_{AC}\) (Moore--Penrose;
\texttt{numerics/optimality/lower\_bound.py}).  For
\((L_A,L_B,L_C)=(4,8,2),(4,12,2),(4,16,2)\) the residual
\(\|Q_{AB}T-Q_{AC}\|\) is \(1.4\times10^{-17},7.0\times10^{-17},2.3\times10^{-17}\),
\(\|T\|=0.552,0.378,0.289<1\), and the two operator inequalities in
\eqref{eq:Tcond} hold with margins \(0.122,0.151,0.163\).  Hence \(E_{AC}=0\)
\emph{exactly} for these geometries, while \(\Erec\approx0.017\,z^2\).  Direct
convex minimisation of the exact \(A\!-\!C\) fidelity over the feasible set gives
\(E_{AC}=1.5\times10^{-9}\) and \(9.8\times10^{-10}\) (solver floor) at
\((4,8,2)\) and \((4,12,2)\), against
\(-\log\Fid(\omega_{AC},\omega_A\otimes\omega_C)=3.2\times10^{-3}\) and
\(1.7\times10^{-3}\): the \(B\)-mediated \(A\!-\!C\) states \emph{contain} the
vacuum \(A\!-\!C\) state.  The route proposed in \texttt{open\_problems\_plan.tex},
Sec.~O8 (``bound the Bures distance between the vacuum \(A\!-\!C\) two-point
function and the set of \(B\)-mediated ones'') therefore yields no bound at all.
The mechanism is transparent: the \(A\!-\!C\) correlations are \(O(z)\) and
\emph{weak}, the \(A\!-\!B\) correlations are \(O(1)\) and \emph{strong}, so
degrading \(B\) into a surrogate \(C\) is easy; what a recovery cannot do is
degrade \(B\) and \emph{simultaneously} keep the \(B\!-\!C\) correlations, and no
single-subalgebra restriction sees that joint constraint.
\end{verdictbreak}

\subsection{Two tools that do survive}

\begin{lemma}[Single-mode measurement bound]\label{lem:onemode}
For any two states \(\varphi,\psi\) on \(\cA(I)\) and any \(h\in\hh_I\),
\(\|h\|=1\), with \(p=\varphi(a^*(h)a(h))\), \(p'=\psi(a^*(h)a(h))\),
\begin{equation}\label{eq:onemode}
 \Fid(\varphi,\psi)\ \le\ \sqrt{pp'}+\sqrt{(1-p)(1-p')},
 \qquad\text{hence}\qquad
 -\log\Fid\ \ge\ \frac{(p-p')^2}{8p(1-p)}+O\bigl((p-p')^3\bigr).
\end{equation}
No gauge invariance, no quasi-freeness and no purity is assumed.
\end{lemma}

\begin{proof}
\lstep{1}{1}{\(n_h:=a^*(h)a(h)\) is a projection.}
\lby{the CAR \(a(h)a^*(h)+a^*(h)a(h)=\|h\|^2=1\) give
\(n_h^2=a^*(h)(1-a^*(h)a(h))a(h)=n_h\) since \(a(h)^2=0\)}
\lstep{1}{2}{\(\{n_h,1-n_h\}\) is a two-outcome POVM with outcome probabilities
\((p,1-p)\) and \((p',1-p')\), and \(\Fid\) is bounded above by the classical
fidelity of the outcome distributions.}
\lby{monotonicity of \(\Fid\) under the unital CP map
\(\mathbb C^2\to\cA(I)\), \((\alpha_0,\alpha_1)\mapsto\alpha_0n_h+\alpha_1(1-n_h)\)
(Lemma~\ref{lem:restriction}, $\langle1\rangle1$); equivalently Fuchs--van~de~Graaf /
Fuchs--Caves}
\lqed{1}{3}
\lby{expanding \(-\log[\sqrt{pp'}+\sqrt{(1-p)(1-p')}]\) at \(p'=p\)}
\end{proof}

\begin{lemma}[Gauge twirling]\label{lem:twirl}
Let \(\gamma_\theta\) be the global \(U(1)\) gauge automorphism, \(\omega\circ\gamma_\theta=\omega\).
For \(\beta\in\mathcal C\) put \(\bar\beta:=\int_0^{2\pi}\frac{d\theta}{2\pi}\,
\gamma_{-\theta}\circ\beta\circ\gamma_\theta\).  Then \(\bar\beta\in\mathcal C\),
\(\bar\beta\) is gauge covariant, and
\(-\log\Fid(\omega,\omega^{\bar\beta})\le-\log\Fid(\omega,\omega^\beta)\).
Hence \eqref{eq:Eopt} may be restricted to gauge-covariant channels.
\end{lemma}

\begin{proof}
\lstep{1}{1}{\(\bar\beta\in\mathcal C\) and is gauge covariant.}
\lby{\(\gamma_\theta\) is a \(*\)-automorphism of both \(\cA(BC)\) and \(\cA(B)\)
(it preserves localisation), so each \(\gamma_{-\theta}\beta\gamma_\theta\) is
admissible; \(\mathcal C\) is convex and closed under the (weak-\(*\) continuous)
average by Prop.~\ref{prop:convex}(1); covariance is the invariance of Haar measure}
\lstep{1}{2}{\(\omega^{\bar\beta}=\int\frac{d\theta}{2\pi}\,\omega^\beta\circ\gamma_\theta\).}
\lby{\(\omega^{\bar\beta}(xy)=\int\frac{d\theta}{2\pi}\omega(x\,\gamma_{-\theta}\beta(\gamma_\theta y))
=\int\frac{d\theta}{2\pi}\omega(\gamma_\theta(x)\,\beta(\gamma_\theta y))
=\int\frac{d\theta}{2\pi}\omega^\beta(\gamma_\theta(x)\gamma_\theta(y))\), using
\(\omega\circ\gamma_{-\theta}=\omega\) and that \(\gamma_\theta\) preserves \(\cA(A)\)}
\lstep{1}{3}{\(\Fid(\omega,\omega^\beta\circ\gamma_\theta)=\Fid(\omega,\omega^\beta)\).}
\lby{invariance of \(\Fid\) under a common automorphism together with
\(\omega\circ\gamma_\theta=\omega\)}
\lqed{1}{4}
\lby{concavity of \(\Fid\) in the second argument (Prop.~\ref{prop:convex}
$\langle2\rangle1$) applied to $\langle1\rangle2$--$\langle1\rangle3$}
\end{proof}

\begin{verdictgap}[\;No reduction to quasi-free channels]
Lemma~\ref{lem:twirl} removes the gauge degrees of freedom but \emph{not} the
non-Gaussian ones.  A ``Gaussification'' step would need a channel-level
analogue of the Wolf--Giedke--Cirac extremality theorem, and none is available:
extremality results say that Gaussian \emph{states} extremise certain
functionals at fixed covariance, not that a Gaussian channel does at least as
well as a general one for the fidelity to a Gaussian target.  Worse, the
obvious structural obstruction is vacuous in the continuum: for a general
gauge-covariant \(\beta\) the recovered \(A\!-\!C\) block is
\(\omega(a^*(f)\,b_g)=\langle a(f)\Omega,\,b_g\Omega\rangle\) with
\(b_g=\beta(a(g))\in\cA(B)\), and by Reeh--Schlieder \(\overline{\cA(B)\Omega}\)
is the whole Hilbert space, so \(b_g\Omega\) is unconstrained except in norm
(Kadison--Schwarz gives only \(\|b_g\Omega\|^2\le\omega^\beta(a^*(g)a(g))\)).
Thus \emph{all} results of Sec.~\ref{sec:numerics} are statements about the
quasi-free class; the lower bound over all channels remains
\(\Erec\ge0\).
\end{verdictgap}



\section{Numerics}\label{sec:numerics}

\subsection{Model, pipeline, validation}

The optimisation needs a model in which the three regions are \emph{exactly}
orthogonal subspaces of the one-particle space; the modular box of Note~5 does
not have this property (its box modes in the coordinate
\(\xi=\log\frac{x+a}{L-x}\) are spread over \(A\), \(B\) and \(C\)), so the
computation is done on the half-filled hopping chain, the model already
validated against the continuum in Note~5 (Sec.~5, agent PL:
\(\texttt{numerics/lattice}\)).  \(A,B,C\) are \(L_A,L_B,L_C\) consecutive sites,
\(Q_{ij}=\sin(\pi(i-j)/2)/(\pi(i-j))\), and \(z=L_AL_C/(L_B(L_A+L_B+L_C))\).
The chain carries both chiralities, so the continuum reference values are
\begin{equation}\label{eq:refvalues}
 \text{Petz}:\ 2\Phi(2z)=8f_2z^2+O(z^3),\qquad
 \text{spatial compression}:\ 2\Phi(z)=2f_2z^2+O(z^3),\qquad
 f_2=\frac{c}{12\pi^2},
\end{equation}
i.e.\ \(8f_2=0.067547\) and \(2f_2=0.016887\) for \(c=1\) (Dirac).  The pipeline
(\texttt{numerics/optimality/}) is:
\begin{enumerate}[leftmargin=1.6em,itemsep=1pt]
\item \texttt{opt\_gaussian.py}: the convex programme of
 Cor.~\ref{cor:convexQP} with capped weights, solved by a log-barrier path
 \(\mu=10^{-1},\dots,10^{-10}\) with damped Newton steps (analytic first
 derivatives, finite-difference Hessian, feasibility backtracking);
\item \texttt{petz\_ref.py}: the Gaussian Petz recovered symbol (Vardhan--Wei--Zou
 eq.~(2.9) at \(\lambda=0\) in the \(T\)-product form of \texttt{petz\_lattice.py})
 and the exact quasi-free root fidelity of Note~5, Thm.~2.1, in \texttt{mpmath}
 at \(90\) digits;
\item \texttt{run\_exact.py}: for each geometry a ladder of caps
 \(\kappa_c\in\{6,\dots,20\}\), each candidate channel evaluated by the exact
 \(-\log\Fid\); the best is reported.  The candidate is a \emph{certificate}: it
 is admissible by Prop.~\ref{prop:cp} and its error is computed exactly, so the
 reported value is a rigorous upper bound for \(\Erec\) in the lattice model.
\end{enumerate}

\paragraph{Validation.}
(i) \emph{Complete positivity and the covariance rule}
(\texttt{cp\_check.py}): for random \((X,Y)\) with \(\dim\hh_B=1\),
\(\dim\hh_D=2\) satisfying \eqref{eq:cp}, the Stinespring dilation of
Prop.~\ref{prop:cp} was built explicitly on a \(4\)-mode Jordan--Wigner Fock
space (\(16\times16\) matrices), the Bogoliubov unitary \(\Gamma(W)=e^{c^*\log W\,c}\)
applied to \(\rho_{AB}\otimes\rho_{Q_E}\), and the two-point function of the
image compared with \eqref{eq:Qrec}: maximal deviation \(6.7\times10^{-16}\).
(ii) \emph{Petz reference}: the recovered symbol reproduces the \(AA\) and the
\(BC\!-\!BC\) blocks of the vacuum symbol to \(0\) (machine exact), as it must,
and \(2\Phi(2z)/(8f_2z^2)=0.933,0.911,0.901\) at \((4,8,2),(6,12,3),(8,16,4)\)
(all at \(z=1/14\)) and \(1.018\) at \((4,16,2)\) (\(z=1/44\)); the continuum
cubic correction alone predicts \(1-0.99\,(2z)=0.858\) at \(z=1/14\), so the
lattice values bracket the continuum law as expected.
(iii) \emph{Second-order formula}: at \((4,6,3)\), \(z=0.154\), the windowed
quantum-Fisher form \eqref{eq:gQ} of the Petz defect at \(\kappa_c=10\) is
\(1.3634\times10^{-3}\) against the exact \(-\log\Fid=1.3032\times10^{-3}\)
(\(+4.6\%\), the sign of the negative cubic remainder \(c_3=-0.99\) of Note~5);
the \emph{uncapped} form gives \(3.61\times10^{-2}\), the ultraviolet
amplification discussed in Sec.~\ref{sec:order}.



\subsection{Results}\label{sec:results}

Table~\ref{tab:opt} collects, for each geometry, the exact \(-\log\Fid\) of the
ordinary Petz map and of the best channel found (over the cap ladder), both
computed at \(90\) digits from the same recovered symbols.  All entries of the
\(\Erec\) column are rigorous \emph{upper} bounds for the quasi-free optimum in
the lattice model.

\begin{table}[h]\centering\small
\begin{tabular}{lccccccc}
\toprule
\((L_A,L_B,L_C)\)&\(n\)&\(z\)&\(-\log\Fid_{\text{Petz}}\)&\(\dfrac{\text{Petz}}{8f_2z^2}\)&
\(-\log\Fid_{\text{best}}\)&\(\dfrac{\text{best}}{2f_2z^2}\)&\(\dfrac{\text{best}}{\text{Petz}}\)\\
\midrule
\((4,8,2)\) &14&\(1/14\)  &\(3.2164\times10^{-4}\)&0.933&\(8.6695\times10^{-5}\)&1.006&\textbf{0.2695}\\
\((4,12,2)\)&18&\(1/27\)  &\(9.1816\times10^{-5}\)&0.991&\(2.6045\times10^{-5}\)&1.124&0.2837\\
\((4,16,2)\)&22&\(1/44\)  &\(3.5524\times10^{-5}\)&1.018&\(1.4684\times10^{-5}\)&1.683&0.4134\\
\((6,12,3)\)&21&\(1/14\)  &\(3.1391\times10^{-4}\)&0.911&\(1.2511\times10^{-4}\)&1.452&0.3986\\
\((6,16,3)\)&25&\(0.045\) &\(1.3069\times10^{-4}\)&0.956&\(5.4347\times10^{-5}\)&1.589&0.4158\\
\((8,16,4)\)&28&\(1/14\)  &\(3.1044\times10^{-4}\)&0.901&\(8.6962\times10^{-5}\)&\textbf{1.009}&\textbf{0.2801}\\
\bottomrule
\end{tabular}
\caption{Exact lattice recovery errors.  \(f_2=c/(12\pi^2)\), \(c=1\);
\(8f_2z^2\) and \(2f_2z^2\) are the continuum two-chirality Petz and spatial
compression values \eqref{eq:refvalues}.  The rows \((6,12,3)\), \((4,16,2)\),
\((6,16,3)\) are limited by the barrier solver at large caps (the objective's
condition number grows like \(e^{\kappa_c}\)); their ``best'' entries are upper
bounds that have not converged (caps reached: \(14,12,12\)).  The three
converged rows give
\(\text{best}/(2f_2z^2)=1.006\), \(1.009\), \(1.124\); the caps attained are
\(\kappa_c=8,10,12\).}
\label{tab:opt}
\end{table}

\paragraph{Reading of the table.}
The reliable rows are those where the cap ladder saturates before the solver
degrades, namely \((4,8,2)\), \((8,16,4)\) and \((4,12,2)\) (the first two at the
same \(z=1/14\) but \(n=14\) and \(n=28\), so the agreement is also a finite-size
check); there
\begin{equation}\label{eq:main-number}
 \boxed{\;\Erec^{\text{lattice}}\;=\;(1.01\text{--}1.12)\times 2f_2z^2,
 \qquad \frac{\Erec}{-\log\Fid_{\text{Petz}}}\;=\;0.270,\ 0.280,\ 0.284
 \;\approx\;\tfrac14 .\;}
\end{equation}
No channel was found anywhere in the search that goes below \(2f_2z^2\); the
best value is \(0.6\%\) above it at \((4,8,2)\).  Since the continuum spatial
geometric compression achieves exactly \(2f_2z^2+O(z^3)\)
(Prop.~\ref{prop:convex}(4)) and the Petz map \(8f_2z^2+O(z^3)\), the numerics
say: \emph{the Petz map is beaten by a factor \(\approx4\), and the winner sits
at the geometric-compression value.}

\paragraph{What the optimal channel is.}
For \((4,8,2)\) the optimiser returns a real \(X\) with singular values
\((1,1,0.9999,0.9998,0.9997,0.9773,0.9768,0.9527)\): \(X\) is, to within a few
percent, an \emph{isometry} \(\hh_B\to\hh_D\), i.e.\ \(X^*\) is a coisometry
\(\hh_D\to\hh_B\) and the channel is very nearly the pure ``squeeze'' of
\(\cA(BC)\) into \(\cA(B)\) with no noise on the eight squeezed directions;
the two remaining directions of \(\hh_D\) (forced by
\(\dim\hh_D-\dim\hh_B=L_C=2\)) receive noise \(Y\le1-XX^*\) with \(\|Y\|=1\),
\(\Tr Y=1.09\).  Writing \(X^*=(S\ \ T)\): \(\|S\|=1\) but
\(\|S-1\|=1.35\) (the channel genuinely moves \(B\) into itself) and
\(\|T\|=0.976\) (the constraint \(T^*T\le1\) is nearly active).  This is exactly
the qualitative signature of the geometric compression: an isometric squeeze
with the range constraint active, \emph{not} a modular rotation (the whole
rotated family collapses onto \(\Phi(\zeta_t)\) with \(\zeta_t\ge z\), Note~5
Cor.~5.1, and never goes below \(\Phi(z)\)).

\paragraph{The Petz property is not optimal.}
The Petz map is characterised by \(\delta_{DD}=Q_{DD}-Z=0\): it reproduces
\(\omega|_{\cA(BC)}\) exactly and puts the entire defect into the \(A\!-\!BC\)
block, \(\|\delta_{AD}\|=9.80\times10^{-3}\) at \((4,8,2)\).  The optimum
\emph{trades}: \(\|\delta_{DD}\|=2.21\times10^{-3}\neq0\) buys
\(\|\delta_{AD}\|=5.62\times10^{-3}\), a \(43\%\) reduction, and
\((5.62/9.80)^2=0.33\) accounts for most of the factor \(0.27\) in the error.
Hence \emph{exact reproduction of the \(BC\) marginal --- the defining property
of the Petz map --- is strictly suboptimal at second order.}

\paragraph{Why the naive lattice compression fails.}
Discretising the continuum M\"obius squeeze \(h_\lambda(x)=x/(1+\lambda x/L)\) in
the site basis (\texttt{compression.py}) produces a matrix whose singular values
run down to \(0.47\)--\(0.62\): on a lattice a squeeze by \((b+c)/b\) pushes
modes through the ultraviolet cutoff, and the resulting channel has
\(-\log\Fid=0.66\)--\(2.2\), three to four orders of magnitude worse than Petz.
The optimum therefore does \emph{not} discretise the geometric map; it realises
the compression's infrared effect while staying isometric in the ultraviolet.
This is the lattice artefact that also explains why the larger systems in
Table~\ref{tab:opt} have not converged: the optimiser must resolve modular modes
up to \(\kappa\sim\kappa_{\max}(n)\), and the barrier Hessian's condition number
grows like \(e^{\kappa_c}\).



\section{Verdicts}\label{sec:verdicts}

\begin{verdictbreak}[\;Q1: is the Petz map second-order optimal? \textbf{No.}]
For the free fermion the ordinary Petz map is beaten by a factor
\(\approx3.5\)--\(3.7\) already in the finite lattice, and by the exact factor
\(4\) in the continuum: the zero-collar geometric compression is admissible
(Prop.~\ref{prop:convex}(4)) and gives \(\Phi(z)=f_2z^2+O(z^3)\) per chirality
against the Petz \(\Phi(2z)=4f_2z^2+O(z^3)\), while the numerics of
Sec.~\ref{sec:results} find no quasi-free channel that does better than the
compression.  Two independent structural reasons: (i) inside the rotated family
the constraint \(\beta(\cA(BC))\subseteq\cA(B)\) is active exactly at the
compression endpoint \(s=\lambda\) (Remark~\ref{rem:active}), and the family is
monotone, so \(t=0\) is the \emph{worst} admissible member for a single
chirality; (ii) the defining property of the Petz map, exact reproduction of the
\(BC\) marginal (\(\delta_{DD}=0\)), is strictly suboptimal: at \((4,8,2)\) the
optimum accepts \(\|\delta_{DD}\|=2.2\times10^{-3}\) in exchange for a \(43\%\)
smaller \(\|\delta_{AD}\|\).  Note that for a \emph{two-dimensional} CFT the
Petz map is still the best member of the \emph{rotated} family (the two
chiralities see opposite rotations, Note~5 eq.~(5.4)); optimality fails only
once one leaves that family.
\end{verdictbreak}

\begin{verdictgap}[\;Q2: is the geometric compression second-order optimal? Consistent with yes, not proved.]
Every channel produced by the convex programme, over six geometries and cap
ladders \(\kappa_c\in\{3,\dots,20\}\), has exact error
\(\ge 2f_2z^2\) (the two-chirality compression value), the closest being
\(1.006\times2f_2z^2\) at \((4,8,2)\) and \(1.124\times\) at \((4,12,2)\);
rows with larger \(n\) are further above only because the barrier solver has not
converged there.  This is evidence, not proof: (a) it is a lattice statement,
(b) it covers the quasi-free class only, (c) the solver's accuracy degrades
exactly where the modular ultraviolet is hardest.  A proof would most naturally
come from the Karush--Kuhn--Tucker conditions of the convex programme of
Cor.~\ref{cor:convexQP} in the continuum modular frame, where the compression is
the point at which the constraint \(1-XX^*\ge0\) is saturated
(\(XX^*=1\), \(Y=0\)) --- i.e.\ one must exhibit a dual certificate supported on
the active constraint.
\end{verdictgap}

\begin{tcolorbox}[colback=navy!5,colframe=navy!60,title={\bfseries Q3: the optimal coefficient},fonttitle=\small]
Per chirality, with \(f_2=c/(12\pi^2)\) and \(z=\eta_{\mathrm V}/(1-\eta_{\mathrm V})\),
\[
 \Erec=\ \kappa_{\mathrm{opt}}\,f_2z^2+O(z^3),\qquad
 \kappa_{\mathrm{opt}}=1.00\ (\pm0.15\ \text{numerically}),\qquad
 \text{versus}\ \kappa_{\text{Petz}}=4 .
\]
For a two-dimensional CFT (both chiralities, spatial channel) the corresponding
statement is \(\Erec=2f_2z^2\) against \(8f_2z^2\).  In the Vardhan--Wei--Zou
variable, \(-\log F_{\mathrm{opt}}=2f_2\eta_{\mathrm V}^2+O(\eta^3)
=0.0169\,c\,\eta_{\mathrm V}^2\) against their \(0.0676\,c\,\eta_{\mathrm V}^2\)
at \(\lambda=0\).  The uncertainty \(\pm0.15\) is the spread of the two converged
rows of Table~\ref{tab:opt} (\(1.006\), \(1.124\)) plus the finite-lattice drift
seen in the Petz column (\(0.90\)--\(1.02\)).
\end{tcolorbox}

\begin{verdictbreak}[\;Q4: lower bound. Only \(\Erec\ge0\) is proved; the planned route is dead.]
The \(A\!-\!C\) restriction proposed in \texttt{open\_problems\_plan.tex}~O8 and
in the Phase-3 brief gives \emph{exactly zero}: the vacuum \(A\!-\!C\) state is
itself \(B\)-mediated (Prop.~\ref{prop:AC-zero}, with the hypothesis verified on
the chain to \(10^{-17}\)).  Gauge twirling (Lemma~\ref{lem:twirl}) is a valid
reduction to gauge-covariant channels but does not reach quasi-free ones, and
the Reeh--Schlieder property makes the naive ``the recovered \(A\!-\!C\)
correlations must factor through \(\hh_B\)'' argument vacuous for non-Gaussian
channels.  The only general tool that survives is the single-mode measurement
bound (Lemma~\ref{lem:onemode}), which converts a lower bound on the recovered
\emph{symbol} error into a lower bound on \(-\log\Fid\) for arbitrary channels;
what is missing is a characterisation of the achievable symbols
\(\{Q^{\mathrm{rec}}(\beta):\beta\in\mathcal C\}\) beyond the quasi-free ones.
\textbf{Open question, sharpened:} is
\(\{Q^{\mathrm{rec}}(\beta):\beta\in\mathcal C\}=\{Q^{\mathrm{rec}}(X,Z):(X,Z)\in\mathcal F\}\)?
If yes, the convex programme of Cor.~\ref{cor:convexQP} plus
Lemma~\ref{lem:onemode} closes O8 completely.
\end{verdictbreak}

\subsection{What is and is not established}

\begin{enumerate}[leftmargin=1.6em,itemsep=2pt]
\item \textbf{Established (proved here).} Convexity of \(\Krec\) and of the
 objective (Prop.~\ref{prop:convex}); the complete-positivity condition
 \eqref{eq:cp} with its Gaussian Stinespring dilation (Prop.~\ref{prop:cp},
 numerically verified); the recovered-symbol rule (Lemma~\ref{lem:rec-symbol});
 the exact convexity of the quasi-free second-order problem in \((X,Z)\)
 (Thm.~\ref{thm:convex-set}) and its semidefinite representability; the
 vanishing of the \(A\!-\!C\) lower bound (Prop.~\ref{prop:AC-zero});
 the single-mode bound (Lemma~\ref{lem:onemode}) and gauge twirling
 (Lemma~\ref{lem:twirl}).
\item \textbf{Numerical (lattice, quasi-free, certificates).} Every entry of the
 \(-\log\Fid_{\text{best}}\) column is an exact evaluation of an admissible
 channel: the factor \(\approx3.7\) improvement over Petz is not a
 solver artefact.
\item \textbf{Not established.} Optimality of the compression (only evidence);
 any nonzero lower bound on \(\Erec\); the extension of the quasi-free
 restriction to all channels; the continuum limit of the lattice optimum (the
 rows at larger \(n\) are unconverged, and the naive continuum-to-lattice
 transcription of the compression fails badly, Sec.~\ref{sec:results}).
\item \textbf{Consequences for the notes.} Note~3's question ``is the Petz map
 optimal?'' is answered in the negative; Note~5's Cor.~5.1 (``the ordinary Petz
 map is worse than the optimal rotation by a factor \(4\)'') is now known to be
 the whole story rather than an artefact of the rotated family, at least among
 quasi-free channels; \texttt{open\_problems\_plan.tex}~O8 should replace the
 \(A\!-\!C\) route by the KKT/dual-certificate route and by the symbol-range
 question above.
\end{enumerate}

\end{document}
